%% notes-tex/eqtl-data-model/sections/1-what-it-is.tex
%% [[torchcell.datamodels.eqtl-data-model]]
%% https://github.com/Mjvolk3/torchcell/tree/main/notes-tex/eqtl-data-model/sections/1-what-it-is.tex
%% SOURCE: hand-authored. Figure from notes/assets/drawio/eqtl-experiment-and-genotype-inference.drawio.
%% Boocock numbers quoted from the mirrored paper.md; schema names read from schema.py.

\section{The experiment and its three tables\secstatus{todo}}\label{sec:what}

Two strains that differ genetically are crossed, classically BY, a laboratory strain
near-isogenic to S288C, against RM, a vineyard isolate
\citep{bremGeneticDissectionTranscriptional2002}. The F1 hybrid is sporulated and
$n$ haploid progeny are collected. Each progeny strain, a \term{segregant}, is a
\term{meiotic mosaic}: at every locus it carries either the BY or the RM allele, in linked
blocks whose boundaries are set by crossover positions. Those boundaries are not known
ahead of time: every segregant draws its own crossover positions in meiosis, so they must
be inferred per strain from the genotype data, and locating them is where the genotype
error concentrates (Sec.~\ref{sec:observed}). Every segregant's transcriptome is
measured, and the question asked per gene is which loci in the mosaic predict its
expression.

The perturbation is therefore not applied to a strain. It is inherited, and it is present
everywhere in the genome at once.

\tcfigfit[1.0]{figures/eqtl-experiment-and-genotype-inference.pdf}{%
  \textbf{An eQTL study, where inference enters it, and what it costs.}
  \textbf{a}, The cross. Both parents are deep-sequenced, so every variant and both its
  alleles are known before any segregant is examined. The segregants are then assayed
  twice, sparsely for genotype and densely for expression, and that asymmetry drives
  everything else.
  \textbf{b}, One segregant along one chromosome, at four levels. True variant sites are
  dense and none is read directly; marker calls are sparse and are read counts rather than
  certainties; the HMM output is a posterior whose ramps are the uncertainty, since a
  crossover is placed only to within one marker interval; the inferred blocks follow. A
  non-crossover gene conversion tract sitting between markers produces no ramp at all,
  because no read reports it.
  \textbf{c}, The two matrices. $X$ shows haplotype blocks rather than noise, so its columns
  are linked; $Y$ is a vector per strain rather than a scalar.
  \textbf{d}, The map that results, and it is the QTL table of
  Sec.~\ref{sec:tables} drawn as a matrix: each dot is one significant (gene, marker) row
  of that table, plotted gene position against marker position. Cis effects fall on the
  diagonal, each variant acting on its own gene and usually through its promoter; trans
  hotspots appear as vertical bands where one regulatory variant moves hundreds of
  transcripts.
  \textbf{e}, The consequence for a record, developed in Sec.~\ref{sec:ontology}.
  Panels \textbf{b} and \textbf{d} are schematic rather than plotted from data. Composed
  from the sources in Sec.~\ref{sec:what} rather than redrawn from a published figure;
  source \file{notes/assets/drawio/eqtl-experiment-and-genotype-inference.drawio}.}
  {fig:eqtl}

\subsection{Two measurements and one inference\secstatus{todo}}\label{sec:tables}

\notesec{Genotype matrix} $X$, segregants by marker positions. Entries encode
\emph{parental origin} rather than base identity, so the natural form looks binary. Bloom
2013 used 11,623 markers for 1,008 segregants \citep{bloomFindingSourcesMissing2013}, and
that older panel is the deliberate reference point rather than a stale one: the
single-cell study of Sec.~\ref{sec:observed} drew its strains from this same cross,
pooling ``393 previously genotyped haploid segregants''
\citep{boocockSinglecellEQTLMapping2025}. BY and RM segregate at roughly
$4.5\times10^{4}$ sites: sites where the two parents carry different alleles, which
therefore separate into different progeny at meiosis (Mendelian segregation) and vary
across the panel. Nothing is swapped or edited at these sites; segregation is inheritance,
not an applied change. The marker set is a thinned representation of that variant
set, and long runs of constant value along the position axis are what recombination blocks
look like. Sec.~\ref{sec:observed} corrects the binary reading: the entries the mapping
models actually consume are posterior probabilities, not calls.

\notesec{Expression matrix} $Y \in \mathbb{R}^{n \times p}$, segregants by genes, with
$p \approx 5{,}000$ to $6{,}000$. This is the phenotype, and it is a vector per strain
rather than a scalar.

\notesec{QTL table} The derived output. For gene $g$ and marker $j$ the standard model is

$$
y_{ig} \;=\; \mu_g \;+\; \beta_{gj}\,X_{ij} \;+\; \varepsilon_{ig},
\qquad \varepsilon_{ig} \sim \mathcal{N}(0,\sigma_g^2)
$$

fit marker by marker, with significance from a LOD score, the $\log_{10}$ odds of the
data under linkage between marker and expression level against the null of no linkage,
and a permutation threshold. A row of the table is a gene, a marker interval, an effect
size, a LOD and a variance explained. Panel d of Fig.~\ref{fig:eqtl} is this table drawn
as a matrix.

The first two are measurements. The third is an inference over them, and it changes if the
mapping method, the threshold or the marker density changes while the measurement does not.
Sec.~\ref{sec:ontology} returns to what that means for storage.

\subsection{Structure worth knowing before modeling it\secstatus{todo}}

\notesec{Cis against trans} A \term{cis} eQTL sits at or near the gene's own locus and is
usually a promoter variant changing that gene's own transcription. A \term{trans} eQTL sits
elsewhere, typically in a regulator. Cis effects are strong, sparse and roughly one per
affected gene; trans effects are weak and numerous.

\notesec{Hotspots} Trans effects are not spread evenly across loci. A small number of loci
each affect hundreds to thousands of transcripts, one regulatory variant propagating
outward. The columns of $Y$ are therefore strongly dependent, a few latent factors explain
much of the variance, and treating genes as independent is wrong. This hotspot structure
is also the leverage genome-scale engineering aims at: a single regulatory locus that
moves hundreds of transcripts.

\subsection{The four candidates, and why they rank where they do\secstatus{todo}}

\begin{table}[H]\centering
\small
\caption[]{The eQTL rows in the candidate list. \emph{Both axes} marks a dataset that is
high-dimensional in perturbation and in readout at once, which in that list is true of
three datasets in total.}
\label{tab:eqtl-rows}
\begin{tabular}{@{}l r l l@{}}
\toprule
Dataset & Segregants & Readout & Both axes \\
\midrule
Albert 2018 \citep{albertGeneticsTransRegulatory2018} & 1,012 & bulk RNA-seq, 1 condition & yes \\
Boocock 2025 \citep{boocockSinglecellEQTLMapping2025} & 393 & one-pot single cell, 2 carbon sources & yes \\
N'Guessan 2025 \citep{nguessanRefiningResolutionYeast2025} & $\sim$4,500 profiled & single cell & yes \\
Smith 2008 \citep{smithGeneEnvironmentInteraction2008} & 109 & bulk, glucose and ethanol & no \\
\bottomrule
\end{tabular}
\end{table}

Boocock 2025 draws 27,744 single cells from those 393 segregants, a median of 17 cells per
segregant. The extra axis is cells \emph{within} a genotype, which is what allows eQTLs to
be mapped for expression noise and for cell-cycle occupancy rather than only for the mean.
Smith 2008 is small, and it is the design that demonstrates an eQTL present in one carbon
source and absent in the other.
